\begin{code}[hide]
module vqupit.sections.symplectic where

open import Data.Nat using (ℕ)
open import Data.Product using (_×_ ; _,_ ; ∃)
open import Data.Sum using (_⊎_)
open import Data.Unit using (⊤ ; tt)
open import Data.Vec using (Vec ; _∷_ ; head ; tail)
open import Relation.Binary.PropositionalEquality using (_≡_)
open import Notations using (₀ ; ₁ ; ₁₊ ; ₂₊)
open import Word.Base using (Word ; ε ; _•_ ; _^_)
open import Presentation.Definitions using (_IsPresentationOf_)

infixl 9 _↑ _↓
infix 8 -_
infixl 6 _+_
infixl 7 _*_
infix 4 _SRel,_===_

private variable
  n : ℕ

data Gate : ℕ → Set where
  H-gate S-gate : Gate 1
  CZ-gate : Gate 2
data Gen : ℕ → Set where
  gate₁ : Gate 1 → Gen (₁₊ n)
  gate₂ : Gate 2 → Gen (₂₊ n)
  _↥   : Gen n → Gen (₁₊ n)

postulate
  proof-elided : ∀ {ℓ} {A : Set ℓ} → A
  A B D E : Set
  Pauli : ℕ → Set
  _+_ _*_ : ∀ {X : Set} → X → X → X
  -_ : ∀ {X : Set} → X → X
  _↑ _↓ : ∀ {n} → Word (Gen n) → Word (Gen (₁₊ n))
  Ex H S CZ Mg M₋₁ : ∀ {n} → Word (Gen n)
  CZ^ S^ Mg^ g^ : ∀ {n} → ℕ → Word (Gen n)
  -b/a p g : ℕ
  act : ∀ {n} → Word (Gen n) → Pauli n → Pauli n
  sform : ∀ {n} → Pauli n → Pauli n → ℕ
  pZ₀ pX₀ : ∀ {n} → Pauli n
  Sp-group : ℕ → Set
  _SRel,_===_ : ℕ → Set

module SimPres where
  postulate presentation : Set
\end{code}
\section{The Symplectic Factor}\label{sec:symplectic}

This section presents the symplectic factor, where the real rewriting
happens, in five steps: the semantics; the normal form; the coset table
that rewrites to it, which is the $S_n$ table of \S\ref{sec:sn} under a
richer action; uniqueness; and the reduction of the relations the table
uses to the eighteen rules that present $\Sp$.

\subsection{The semantics: symplectic transformations of Pauli vectors}\label{sec:spsem}

Rather than $2n \times 2n$ matrices, the formalisation models $\Sp$ as
the group of $\Zp$-linear, symplectic-form-preserving bijections of the
phase space $\aid{Pauli}\;n = \aid{Vec}\;(\Zp \times \Zp)\;n$ --- one
pair $(a,b)$, the exponents of $Z^{a}X^{b}$ (up to phase), per wire.
A symplectic transformation is a record \aid{Symplectic}~$n$ carrying
the action \aid{ap}, its inverse, and proofs that the two are inverse,
that \aid{ap} is linear, and that it preserves the symplectic form
\aid{sform}.  Equality is extensional equality of the action alone, so
the proof fields never obstruct.  Composition is function composition,
and therefore associativity and the unit laws hold \emph{on the nose},
by \aid{refl} --- a small choice with large consequences downstream,
since every monoid case of every soundness proof closes definitionally.
\aid{Sp-group}~$n$ is the resulting \aid{Group}.  The generators act
as the qupit paper's conjugation tables (its Lemmas~2.22--2.23) read on
exponent vectors, and $\sem{\_}$ is the homomorphic extension:

\begingroup\renewcommand{\AgdaCodeStyle}{\footnotesize}
\begin{code}
actg : ∀ {n} → Gen n → Pauli n → Pauli n
actg (gate₁ H-gate)  ((a , b) ∷ ps)             = (- b , a) ∷ ps
actg (gate₁ S-gate)  ((a , b) ∷ ps)             = (a , b + a) ∷ ps
actg (gate₂ CZ-gate) ((a , b) ∷ (a' , b') ∷ ps) = (a , b + a') ∷ (a' , b' + a) ∷ ps
actg (g ↥)           (x ∷ ps)                   = x ∷ actg g ps
\end{code}
\endgroup

\noindent Words are read left to right, so $w \bullet v$ applies $w$
first and its denotation composes on the right.  The seven proof
obligations per
generator --- inverses, linearity, preservation of the form --- are
identities in $\Zp$, discharged with the ring solver.  Soundness of
every axiom of \S\ref{sec:reduction} in this semantics is one clause per
axiom in \texttt{SoundnessDirect} (1\,100 lines).

\subsection{The normal form: boxes, rows, and the staircase}\label{sec:nf}

The qupit paper's symplectic normal form is a $Z$-normal layer followed
by an $X$-normal layer on $n$ wires, then a normal form on $n-1$ wires.
We refine this into a doubly inductive datatype.  The atomic boxes were
shown in \S\ref{sec:background}; rows and the staircase are:

\begin{minipage}[c]{0.4\textwidth}
\begingroup\renewcommand{\AgdaCodeStyle}{\footnotesize}
\begin{code}
M L ML : ℕ → Set
M 0 = ⊤
M (₁₊ n) = Vec D n × E
L 0 = ⊤
L (₁₊ n) = Vec B n × A
ML 0 = ⊤
ML 1 = M 1 × L 1
ML m@(₂₊ n) = M m × L m
            ⊎ D × ML (₁₊ n)
\end{code}
\endgroup
\begin{code}[hide]
postulate
  [_]ᵐˡ : ∀ {n} → ML n → Word (Gen n)
\end{code}
\end{minipage}\hfill
\begin{minipage}[c]{0.56\textwidth}\centering
  \begin{tabular}{@{}r@{\,}c@{\,}l@{\quad}r@{\,}c@{\,}l@{}}
    \ufig[0.5]{m4-box} & \scalebox{0.7}{$=$} & \ufig[0.5]{m4-atoms} &
    \ufig[0.5]{ml4-box} & \scalebox{0.7}{$=$} & \ufig[0.5]{ml-tower-a-r}\\[1ex]
    \ufig[0.5]{l4-box} & \scalebox{0.7}{$=$} & \ufig[0.5]{l4-atoms} &
    \ufig[0.5]{ml4-box} & \scalebox{0.7}{$=$} & \ufig[0.5]{ml-tower-b-r}\\
  \end{tabular}
\end{minipage}

\medskip\noindent
An $M$-row is a column of $D$ boxes ending in an $E$ box (the paper's
$X$-normal layer); an $L$-row is a column of $B$ boxes ending in an $A$
box (its $Z$-normal layer, with the $A$ box at the bottom).  An
$\ML$ box at width $m \ge 2$ is \emph{either} a full pair of rows
\emph{or} a bottom $D$ box followed by an $\ML$ box one wire up: a
spine of $D$ boxes capped by a pair of rows.  This is the first
induction, on the width of a single coset representative.  The second
induction is the staircase:

\begin{minipage}[c]{0.54\textwidth}
\begingroup\renewcommand{\AgdaCodeStyle}{\footnotesize}
\begin{code}
NF : ℕ → Set
NF 0 = ⊤
NF (₁₊ n) = NF n × ML (₁₊ n)

[_] : ∀ {n} → NF n → Word (Gen n)
[_] {0}     tt         = ε
[_] {₁₊ n}  (nf , ml)  = [ nf ] ↑ • [ ml ]ᵐˡ
\end{code}
\endgroup
\end{minipage}\hfill
\begin{minipage}[c]{0.42\textwidth}\centering
  \ufig[0.7]{nf4}
\end{minipage}

\medskip\noindent
Reading normal forms back as circuits is structural recursion, box by
box.  The box-to-word maps make the pictures of \S\ref{sec:background}
precise; for instance
\begin{code}
[_]ᵈ : ∀ {n} → D → Word (Gen (₂₊ n))
[_]ᵈ {n} (₀ , b) = Ex • CZ^ (- b)
[_]ᵈ {n} (a@(₁₊ _) , b) = Ex • CZ^ (- a) • H • S^ -b/a
\end{code}
where $-b/a = -b \cdot a^{-1}$, and an $A$ box $A_{a,b}$ with $a \ne 0$
is $\aid{XM}\;a \bullet H \bullet S^{-b/a}$, a multiplier, a Hadamard,
and an $S$-power.  Why the spine?  The paper's $Z$-normal circuit
$W_Z^{(n)}$ has its $A$ box on some wire $i$, with $B$ boxes above it
and nothing below; the position $i$ is an index that the paper's
normal-form circuits carry implicitly.  The $\ML$ type makes it explicit
as the length of the $D$-spine, and, crucially, makes the coset type at
width $m$ a \emph{sum} of the coset type at width $m-1$ and a flat case
--- exactly the shape the coset table and the uniqueness proof recurse
on.  The identity coset is the flat one with all exponents zero and
$A_{0,1}$ at the bottom, and the section sends it to $\varepsilon$ up to
$\approx$, because $\mathit{Ex}^2 \approx \varepsilon$ collapses each
$D_{0,0}$ and $B_{0,0}$.

\subsection{From symmetric to symplectic: change the action}\label{sec:pushing}

The coset table of \S\ref{sec:sn} serves the symplectic group under a
richer action.  There are three ways of looking at the wires.  $S_n$
acts on $\{0,\dots,n-1\}$: a wire is a position, and $c \circ f^{-1}$
fixes $0$.  $S_n$ acts on $\aid{Vec}\;\mathsf{Bit}\;n$: a wire is a
bit, and $c \circ f^{-1}$ fixes the $\mathbb{Z}_2$-linear subspace of
the bottom wire.  Adding $H$, $S$, $CZ$, the group acts on
$\aid{Vec}\;\mathrm{Pauli}_1\;n$, a $2n$-dimensional symplectic
$\Zp$-vector space in which a wire is a two-dimensional symplectic
subspace; enlarging the coset layer from staircases to $\ML$ boxes, for
each $f \in \Sp$ there is a $c$ such that $c \circ f^{-1}$ fixes the
symplectic subspace of the bottom wire.  The cosets at width $1{+}n$
are the $\ML$ boxes, their section is $[\_]^{\mathsf{ml}}$, and the coset
table is the paper's ``push a gate through a normal box'', now a
function $\aid{ract} : \ML\,(1{+}n) \to \aid{Gen}\,(1{+}n) \to
\aid{Circuit}\;n \times \ML\,(1{+}n)$ with the same soundness law
$[c]^{\mathsf c} \bullet g \approx b'{\uparrow} \bullet [c']^{\mathsf
c}$ as in \S\ref{sec:sn}.  \aid{ract} dispatches on the shape of the coset --- a flat box or a
$D$-spine --- and on the generator --- $S$ or $H$ on the bottom wire,
$CZ$ on the bottom two, or a lifted gate --- and delegates each case to
a module of the \texttt{Pushing} directory (26 files, 5\,600 lines):
$S$ and $H$ through a flat box, $CZ$ through a flat box, $CZ$ meeting a
$D$ box and a flat box, a lifted gate through the spine.  The residual
$b'$ lives one wire up, as the type says.  The soundness proof
\aid{ract-sound} is a chain of $\approx$-steps per clause, each step a
box relation or a structural rule.  For example, an $S$ on the bottom
wire meeting a spine $D_{a,b}$ followed by $\mathit{lm}$: $S$ commutes
past the lifted $\mathit{lm}$, then the box relation
$D \leftarrow S$ rewrites $D_{a,b} \bullet S$ to $\mathit{dir}{\uparrow}
\bullet D_{a,b-a}$, and re-association finishes.

So far this is the $S_n$ story with bigger cosets.  The difference is in
the five hypotheses of \S\ref{sec:engine}.  Four close as before; the
second, \aid{h-wd-ax} --- that the table respects the axioms, i.e.\
congruent inputs produce equal cosets and congruent residuals --- does
\emph{not} close by evaluation, because the axioms of $\Sp$ (unlike
$\sigma^2 = e$) are not decided by a finite computation on the coset
data.  The formalisation proves it in two halves.  The \emph{coset half}
is proved semantically: apply $\sem{\_}$ to \aid{ract-sound} on both
sides of an axiom, observe that residuals live one wire up and so cannot
move the bottom wire, and conclude from head-injectivity of the coset
section (\S\ref{sec:uniqueness}) that the two cosets are equal.  The
\emph{residual half} then follows by cancellation --- provided the shift
$\uparrow$ is injective on congruence classes at the lower width, which
is a consequence of faithfulness (completeness) there.  This makes the
tower a well-founded induction on the width, \aid{tower : ∀ n →
TowerLevel n}, where a level packages the normal form at that width
together with the agreement of its section with the coset sections: a
normal form at width $k$ gives faithfulness at $k$ (by
\aid{by-normalization}); faithfulness gives injectivity of $\uparrow$
on congruence classes (using that a lifted word acts as $\mathrm{id}
\oplus \sem{w}$); that discharges \aid{h-wd-ax} at width $k$; and the
generic normal-form transport produces the normal form at width
$k{+}1$.  The base case, width $1$, is built by hand.  This is the
formal counterpart
of the paper's Lemma~4.2 and Proposition~4.3: no termination argument is
needed beyond structural recursion, and ``the dirty gates eventually get
removed'' is the type of \aid{ract}.

\subsection{Uniqueness of the normal form}\label{sec:uniqueness}

Theorem~\ref{thm:unique} says that normal forms with equal denotations
are equal.  The action of a normal form on a Pauli vector peels the
staircase one wire at a time --- the coset factor acts, the resulting
bottom entry is fixed, and the inner normal form acts on the tail ---
and injectivity is a structural induction on the staircase:

\begingroup\renewcommand{\AgdaCodeStyle}{\footnotesize}
\begin{code}
lemma-lm-head-inj : ∀ {n} (lm₁ lm₂ : ML (₁₊ n)) →
  (∀ (ps : Pauli (₁₊ n)) → head (act [ lm₁ ]ᵐˡ ps) ≡ head (act [ lm₂ ]ᵐˡ ps)) → lm₁ ≡ lm₂
lemma-lm-tail-surj : ∀ {n} (lm : ML (₁₊ n)) (qs : Pauli n) →
  ∃ λ (ps : Pauli (₁₊ n)) → tail (act [ lm ]ᵐˡ ps) ≡ qs
\end{code}
\endgroup
\begin{code}[hide]
lemma-lm-head-inj = proof-elided
lemma-lm-tail-surj = proof-elided
\end{code}

Two obligations.  \emph{Tail surjectivity} is free, because a
symplectic transformation carries its inverse.  \emph{Head injectivity}
--- the bottom-wire output determines the coset --- is the crux, and it
is itself an induction on the $\ML$ type (\texttt{LMHeadInj},
1\,300 lines): at width $1$ a case analysis on the four shapes of an
$A$ box using closed forms for the actions of $S^{k}$, $HS^{k}$ and the
multipliers; a separation lemma showing a flat box and a spine never
agree on heads; and, for two spines, the $D$ box is read off by probing
with the two basis Paulis $Z$ and $X$ on the bottom wire and the identity
above, after which the hypothesis descends one width.  With these two
lemmas the generic \texttt{Circuit.Uniqueness} module of the library
assembles \aid{⟦[]⟧-injective} and the \aid{UniqueNormalForm} witness;
the module's header records the one subtlety, that \aid{ap} is a
\emph{left} action --- \aid{ap} of a composite applies its right factor
first --- so the coset factor of a normal form is the one next to the
input, which is what lets the coset be recovered from every state rather
than from states of a restricted shape.  This is the paper's
Proposition~3.8 with its Lemmas~3.4--3.7, now an induction whose
structure is dictated by the datatype.

\subsection{Relation reduction: box relations and the two rule sets}\label{sec:reduction}

Which relations does \aid{ract-sound} use?  Rewriting \emph{every}
circuit to normal form and collecting every relation the rewrite ever
uses gives the paper's box relations --- one per gate per box shape,
parametrised by the box's labels.  The formalisation organises them as
nine families of case tables, each a lemma quantified over the box
parameters with a computed direction \aid{dir-of} and updated box:

\begin{center}\small
\begin{tabular}{@{}lll@{}}
\toprule
wires & family & pushed \\
\midrule
1 & $A \leftarrow H$, $A \leftarrow S$, $E \leftarrow S$ & $H$ or $S$ through an $A$ box; $S$ through an $E$ box \\
2 & $B \leftarrow H{\uparrow},S{\uparrow},S$; \ $D \leftarrow H,S{\uparrow},S,CZ$; \ $L \leftarrow CZ$ & a gate through a $B$, a $D$, or an $L$-row \\
3 & $BB \leftarrow CZ{\uparrow}$, $B{\uparrow} \leftarrow CZ$, $DD \leftarrow CZ$ & $CZ$ through two $B$ or two $D$ boxes \\
\bottomrule
\end{tabular}
\end{center}

For instance, the two-wire table for a $D$ box states that $[d]^{\mathsf
d} \bullet g \approx S^{e} \bullet \mathit{dir}{\uparrow} \bullet
[d']^{\mathsf d}$ with $e$, $\mathit{dir}$ and $d'$ computed from the
labels of $d$ and the gate $g$ --- the paper's Figure~18, with the
residual split into an $S$-power on the bottom wire and a word one wire
up.  The nine families are the paper's
Figures~15--21; unfolding the case tables gives 42 parametric relations,
against the paper's count of 66, the difference being only in how cases
are grouped (the formalisation, for instance, proves each of the
$\uparrow$/$\downarrow$ pairs of a $B$ box and a $D$ box once and
obtains the other by a duality lemma).  The width-3 relations are stated
at width~3 and widened to arbitrary width by a congruence
$\aid{cong↓ᵏ}$.

The box relations are derived (\texttt{BR}, 7\,000 lines, on top of a
calculus of $\mathit{Ex}$-conjugations and two-qupit lemmas of 11\,000
lines) from the \emph{full} symplectic rule set, seventeen group-specific
axioms whose two- and three-qupit members are Selinger's
c10--c15~\cite{selinger2015clifford} verbatim and whose one-qupit
members are the multiplier calculus --- $S^{p} = H^{4} = (SH)^{3} =
\varepsilon$, $H^{2}S = SH^{2}$, and four \emph{schemes} over units $x,
y$: $M_x M_y = M_{xy}$, $M_x S = S^{x^2} M_x$, and $M_x$ on either wire
past $CZ$ gives $CZ^{x}$ --- with $M_x$ spelled $S^{x}HS^{x^{-1}}HS^{x}H$.
By inspection one finds a smaller set of primitive, bidirectional
relations that derive all of these: the \emph{simplified} symplectic
rule set, which replaces the schemes by statements about the single
multiplier $M_g$ of the primitive root, strengthens $H^4 = \varepsilon$
to $H^2 = M_{-1}$, and drops $(SH)^{3} = \varepsilon$ and $H^2S =
SH^2$:

\begin{code}
order-S :           (₁₊ n) SRel,  S ^ p === ε
order-H :           (₁₊ n) SRel,  H ^ 2 === M₋₁
M-power :     ∀ k → (₁₊ n) SRel,  Mg^ k === M (g^ k)
semi-MS :           (₁₊ n) SRel,  Mg • S === S^ (g * g) • Mg
semi-M↑CZ :         (₂₊ n) SRel,  Mg ↑ • CZ === CZ^ g • Mg ↑
semi-M↓CZ :         (₂₊ n) SRel,  Mg ↓ • CZ === CZ^ g • Mg ↓
\end{code}
\begin{code}[hide]
order-S = proof-elided
order-H = proof-elided
M-power = proof-elided
semi-MS = proof-elided
semi-M↑CZ = proof-elided
semi-M↓CZ = proof-elided
\end{code}

plus the same \aid{order-CZ}, \aid{comm-CZ-S↓}, \aid{comm-CZ-S↑} and
\aid{selinger-c10} to \aid{c15}: fifteen group-specific rules, or
eighteen with the three structural ones.  The choice is not canonical;
one picks the set one finds natural.  The two rule sets generate the
same congruence, and since they share an alphabet the identity on words
is an isomorphism of presentations in both directions
(\texttt{Simplified.Iso}); the substantive direction derives each full
axiom from the simplified ones.  The multiplier law $M_x \bullet M_y
\approx M_{xy}$ is typical: write $x = g^{k}$, $y = g^{l}$, use
\aid{M-power} twice, add the exponents, reduce modulo $p-1$ using
$M_g^{\,p-1} \approx \varepsilon$, which is Fermat's little theorem, and
use \aid{M-power} once more.  The whole reduction from seventeen rules
to fifteen is 900 lines; the reduction from the box relations to the
seventeen is 18\,000.  This is the formal counterpart of the paper's
Theorem~4.4, whose proof the paper already delegated to Agda.  Read
modulo Paulis and scalars, Selinger's three-qubit
relations~\cite{selinger2015clifford} are \emph{identical} to
\aid{selinger-c12} to \aid{c15}, which is why the names were kept, and
his two-qubit relations are similar; the one-qupit relations are where
the odd prime shows, in the multiplier calculus of $M_g$ with its
dependence on a primitive root and on Fermat's theorem, which has no
qubit counterpart.  That the difficulty of the general odd prime
concentrates on one wire is what makes the factored route affordable:
the two- and three-wire rewriting is Selinger's, and the one-wire
arithmetic is arithmetic.

\subsection{The symplectic presentation theorem}\label{sec:sptheorem}

Soundness (one clause per axiom), the normal form of \S\ref{sec:pushing}
and its uniqueness (\S\ref{sec:uniqueness}) assemble, through
\aid{by-normalization}, into a sub-presentation: the interpretation is an
injective homomorphism from the words modulo the full rules into
\aid{Sp-group}.  Surjectivity --- every symplectic transformation is the
action of some circuit --- is proved constructively by ``clear the bottom
wire, recurse on the tail'': the single-level input is

\begin{code}
Theorem-LM : ∀ {n} (p q : Pauli n) → sform p q ≡ ₁ →
  ∃ λ (lm : ML n) → act [ lm ]ᵐˡ p ≡ pZ₀ × act [ lm ]ᵐˡ q ≡ pX₀
\end{code}
\begin{code}[hide]
Theorem-LM = proof-elided
\end{code}

--- for any symplectic pair there is a coset box sending it to the
bottom-wire $Z$ and $X$ --- which is the paper's Lemma~3.7 with the box
computed rather than asserted.  Composing the sub-presentation of the
full rules with the identity isomorphism of \S\ref{sec:reduction} gives
the theorem for the simplified rules, with the \emph{same}
interpretation:

\begin{code}
Theorem-Sp : ∀ {n} → (n SRel,_===_) IsPresentationOf (Sp-group n)
Theorem-Sp = SimPres.presentation
\end{code}

